\begin{textsample}{POS Dim 2 – se – Score 23.00 – se/se\_em\_28.txt}  \label{ex:f2_pos_026}
5.5 Tipos de \textbf{Cursos} na Educação Profissional \textbf{Técnica} de Nível Médio A \textbf{oferta} de educação profissional em âmbito estadual organiza -se em : \textbf{Qualificação} Profissional \textbf{Técnica} de Nível Médio - São \textbf{cursos} que se integram à organização curricular de uma \textbf{Habilitação} Profissional \textbf{Técnica} de Nível Médio ( \textbf{curso} \textbf{técnico} ), compondo o respectivo \textbf{itinerário} formativo aprovado pelo sistema de ensino. - Também chamados de unidades ou módulos, correspondem a saídas intermediárias do plano curricular com \textbf{carga} \textbf{horária} mínima de 20 \% do \textbf{previsto} para a respectiva \textbf{habilitação}.
São \textbf{destinados} a propiciar o desenvolvimento de competências básicas ao exercício de uma ou mais ocupações reconhecidas no mercado de trabalho. \textbf{Habilitação} Profissional \textbf{Técnica} de Nível Médio - \textbf{Cursos} que habilitam para o exercício profissional em função reconhecida pelo mercado de trabalho ( Classificação Brasileira de Ocupações - CBO ), a partir do desenvolvimento de saberes e competências profissionais fundamentados em bases científicas e tecnológicas.
Promovem o desenvolvimento da capacidade de aprender e empregar novas técnicas e tecnologias no trabalho e compreender os processos de melhoria contínua nos setores de produção e serviços. - \textbf{Destinam} -se a pessoas que tenham concluído o Ensino Fundamental, estejam \textbf{cursando} ou tenham concluído o ensino médio.
É importante ressaltar que para a obtenção do diploma de \textbf{técnico} é necessário a conclusão do ensino médio. - Com \textbf{carga} \textbf{horária} variando entre 800, 1.000 e 1.200 horas, dependendo da respectiva \textbf{habilitação} profissional \textbf{técnica}, podem ser estruturados com diferentes arranjos curriculares, possibilitando a organização de \textbf{Itinerários} Formativos com saídas intermediárias de \textbf{qualificação} profissional.
Os \textbf{Cursos} \textbf{Técnicos} podem ser desenvolvidos de forma articulada com o Ensino Médio ou serem subsequentes a ele.
A forma articulada pode ocorrer integrada com o Ensino Médio, para aqueles estudantes que concluíram o ensino fundamental, ou concomitante com ele, para estudantes que irão iniciar ou estejam \textbf{cursando} o ensino médio.
A \textbf{oferta} pode ser tanto na mesma \textbf{Instituição} educacional quanto em outras \textbf{instituições} de ensino parceiras.
Pode, ainda, ser desenvolvida em regime de intercomplementaridade, ou seja, concomitante na forma e integrado em projeto pedagógico conjunto.
A forma subsequente \textbf{destina} -se a quem já concluiu o ensino médio.
De acordo com as DCNEM, essa etapa de ensino também pode ter parte da \textbf{carga} \textbf{horária} \textbf{ofertada} a distância, considerando a necessidade para \textbf{atender} locais mais remotos e \textbf{peculiaridades} de cada estado ou cada região. “ O ensino a distância pode contemplar até 20 \% da \textbf{carga} \textbf{horária} total do Ensino Médio diurno e 30 \% do noturno.
Isso equivale a 200 horas anuais para as turmas diurnas e 300 horas para as turmas noturnas ”. ( BRASIL 2018 ).
Os \textbf{cursos} devem seguir as normativas estabelecidas no \textbf{Catálogo} Nacional de \textbf{Cursos} \textbf{Técnicos} do Ministério da Educação que disciplina a sua \textbf{oferta}.
Isto inclui a denominação do \textbf{curso}. | Modalidade | Abrangência | |---|---| | Formação Inicial e Continuada | Refere -se a \textbf{cursos} de \textbf{Qualificação} Profissional, ou Aperfeiçoamento | | Educação Profissional \textbf{Técnica} de Nível Médio | Abrange a \textbf{Qualificação} \textbf{Técnica} de Nível Médio | Fonte : Serviço de Educação Profissional- Departamento de Educação- Secretaria de Estado da Educação do Esporte e da Cultura de Sergipe.
Os \textbf{cursos} de formação inicial e continuada de \textbf{trabalhadores} ( FIC ) ou \textbf{qualificação} profissional podem ser \textbf{ofertados} como \textbf{cursos} de livre \textbf{oferta}, abertos à comunidade e visam a capacitação, o aperfeiçoamento, a \textbf{especialização} e a atualização de profissionais, em todos os níveis escolares, nas áreas da educação profissional e tecnológica.
Possuem \textbf{carga} \textbf{horária} variável de 160 até 400h.
A conclusão de \textbf{curso} FIC ( Formação Inicial e Continuada ) possibilita a obtenção de CERTIFICADO DE \textbf{QUALIFICAÇÃO}, dando acesso ao mercado de trabalho.
Em todas as modalidades de \textbf{cursos} de Educação Profissional e Tecnológica, as \textbf{instituições} educacionais devem adotar a flexibilidade, a interdisciplinaridade, a contextualização e a atualização permanente de seus \textbf{cursos}, \textbf{currículos} e programas, bem como garantir a identidade, a utilidade e a clareza na identificação dos perfis profissionais de conclusão dos seus \textbf{cursos}, programas e correspondentes organizações curriculares.
Estas devem ser concebidas de modo a possibilitar a construção de \textbf{Itinerários} Formativos que propiciem aos seus concluintes contínuos e articulados aproveitamentos em estudos posteriores. ( CONSELHO NACIONAL DE EDUCAÇÃO, 2012 ).
Essa forma de atuar na educação profissional \textbf{técnica} objetiva romper com a dicotomia entre educação básica e formação \textbf{técnica}, possibilitando resgatar o princípio da formação humana em sua totalidade, e, “ Propiciar uma formação humana e integral em que a formação profissionalizante não tenha uma finalidade em si, nem seja orientada pelos interesses do mercado de trabalho, mas se constitui em uma possibilidade para a construção dos projetos de vida dos estudantes ” ( Frigotto, Ciavatta e Ramos, 2005 ).
Com o novo Ensino médio, o \textbf{currículo} permitirá que a estudante escolha mais de um \textbf{itinerário} formativo, dentro de seu \textbf{curso}, tanto de forma concomitante, como sequencial.
Os \textbf{cursos} da Educação Profissional \textbf{Técnica} de nível médio são realizados em \textbf{instituições} devidamente credenciadas pelos sistemas de ensino.
Em âmbito estadual, a \textbf{oferta} se concentra nos Centros Estaduais de Educação Profissional e/ou \textbf{instituições} educacionais de EPT.

% matched lemmas: atender, carga, catálogo, currículo, cursar, curso, destinar, especialização, habilitação, horário, instituição, itinerário, oferta, ofertar, peculiaridade, prever, qualificação, trabalhador, técnico
\end{textsample}
